\documentclass[10pt,a4paper]{amsart}
\usepackage{amsmath,amssymb,amsthm,graphicx}
\usepackage{verbatim}
\newenvironment{reference}{\comment}{\endcomment}
%\newenvironment{reference}{}{}
\newenvironment{slogan}{\comment}{\endcomment}
\newenvironment{history}{\comment}{\endcomment}

% For commutative diagrams we use Xy-pic
\usepackage[all]{xy}

% We use 2cell for 2-commutative diagrams.
\xyoption{2cell}
% Theorem environments.
%
\theoremstyle{plain}
\newtheorem{theorem}[subsection]{Theorem}
\newtheorem{proposition}[subsection]{Proposition}
\newtheorem{lemma}[subsection]{Lemma}

\theoremstyle{definition}
\newtheorem{definition}[subsection]{Definition}
\newtheorem{example}[subsection]{Example}
\newtheorem{exercise}[subsection]{Exercise}
\newtheorem{situation}[subsection]{Situation}

\theoremstyle{remark}
\newtheorem{remark}[subsection]{Remark}
\newtheorem{remarks}[subsection]{Remarks}

\numberwithin{equation}{subsection}

% Macros
%
\def\lim{\mathop{\mathrm{lim}}\nolimits}
\def\colim{\mathop{\mathrm{colim}}\nolimits}
\def\Spec{\mathop{\mathrm{Spec}}}
\def\Hom{\mathop{\mathrm{Hom}}\nolimits}
\def\Ext{\mathop{\mathrm{Ext}}\nolimits}
\def\SheafHom{\mathop{\mathcal{H}\!\mathit{om}}\nolimits}
\def\SheafExt{\mathop{\mathcal{E}\!\mathit{xt}}\nolimits}
\def\Sch{\mathit{Sch}}
\def\Mor{\mathop{\mathrm{Mor}}\nolimits}
\def\Ob{\mathop{\mathrm{Ob}}\nolimits}
\def\Sh{\mathop{\mathit{Sh}}\nolimits}
\def\NL{\mathop{N\!L}\nolimits}
\def\CH{\mathop{\mathrm{CH}}\nolimits}
\def\proetale{{pro\text{-}\acute{e}tale}}
\def\etale{{\acute{e}tale}}
\def\QCoh{\mathit{QCoh}}
\def\Ker{\mathop{\mathrm{Ker}}}
\def\Im{\mathop{\mathrm{Im}}}
\def\Coker{\mathop{\mathrm{Coker}}}
\def\Coim{\mathop{\mathrm{Coim}}}

% Boxtimes
%
\DeclareMathSymbol{\boxtimes}{\mathbin}{AMSa}{"02}

%
% Macros for moduli stacks/spaces
%
\def\QCohstack{\mathcal{QC}\!\mathit{oh}}
\def\Cohstack{\mathcal{C}\!\mathit{oh}}
\def\Spacesstack{\mathcal{S}\!\mathit{paces}}
\def\Quotfunctor{\mathrm{Quot}}
\def\Hilbfunctor{\mathrm{Hilb}}
\def\Curvesstack{\mathcal{C}\!\mathit{urves}}
\def\Polarizedstack{\mathcal{P}\!\mathit{olarized}}
\def\Complexesstack{\mathcal{C}\!\mathit{omplexes}}
% \Pic is the operator that assigns to X its picard group, usage \Pic(X)
% \Picardstack_{X/B} denotes the Picard stack of X over B
% \Picardfunctor_{X/B} denotes the Picard functor of X over B
\def\Pic{\mathop{\mathrm{Pic}}\nolimits}
\def\Picardstack{\mathcal{P}\!\mathit{ic}}
\def\Picardfunctor{\mathrm{Pic}}
\def\Deformationcategory{\mathcal{D}\!\mathit{ef}}
\usepackage{cleveref}
\crefname{originallemma}{original lemma}{original lemmas}
\Crefname{originallemma}{Original lemma}{Original lemmas}
\makeatletter
\newlabel{coherent-lemma-quasi-coherent-affine-cohomology-zero}{{01XB}{source}{}{sourcepointer-1}{}}
\newlabel{coherent-lemma-quasi-coherent-affine-cohomology-zero@cref}{{[originallemma][0][]01XB}{[1][0][]source}{}{}{}}
\makeatother
\crefname{originallemma}{original lemma}{original lemmas}
\Crefname{originallemma}{Original lemma}{Original lemmas}
\makeatletter
\newlabel{sites-cohomology-lemma-derived-limit-suitable-system}{{0A09}{source}{}{sourcepointer-2}{}}
\newlabel{sites-cohomology-lemma-derived-limit-suitable-system@cref}{{[originallemma][0][]0A09}{[1][0][]source}{}{}{}}
\makeatother
\crefname{originallemma}{original lemma}{original lemmas}
\Crefname{originallemma}{Original lemma}{Original lemmas}
\makeatletter
\newlabel{sites-cohomology-lemma-inverse-limit-is-derived-limit}{{0BKY}{source}{}{sourcepointer-3}{}}
\newlabel{sites-cohomology-lemma-inverse-limit-is-derived-limit@cref}{{[originallemma][0][]0BKY}{[1][0][]source}{}{}{}}
\makeatother
\crefname{originalsection}{original section}{original sections}
\Crefname{originalsection}{Original section}{Original sections}
\makeatletter
\newlabel{spaces-cohomology-section-higher-direct-image}{{071Y}{source}{}{sourcepointer-4}{}}
\newlabel{spaces-cohomology-section-higher-direct-image@cref}{{[originalsection][0][]071Y}{[1][0][]source}{}{}{}}
\makeatother
\title{Derived limits of surjective inverse systems of quasi-coherent complexes}
\author{The Stacks Project Authors}
\date{stacks-project:spaces-perfect.tex:local-snapshot-20261012; GNU FDL 1.2 or later}
\begin{document}
\maketitle
\noindent\textit{Editorial note.} All delivered mathematical text is the author's own, copied from the pinned Stacks Project chapter file spaces-perfect.tex; no statement or proof step was rewritten, reordered or removed.\par\medskip
\noindent\textit{Editorial note.} The counted proof is the author's complete argument: the proof verifies on affine objects the vanishing of higher cohomology of quasi-coherent sheaves and the surjectivity of sections under surjective maps, then applies the two Cohomology on Sites lemmas for the three assertions\par\medskip
\noindent\textit{Editorial note.} The label lemma-Rlim-quasi-coherent names the counted statement, and the proof environment immediately following it in the source is delivered verbatim as its proof.\par\medskip
\noindent\textit{Editorial note.} The counted statement is carried by the author's own lemma environment in the chapter source.\par\medskip
\noindent\textit{Allowed standard prerequisites (editorial).}\begin{itemize}
\item vanishing of higher cohomology of quasi-coherent sheaves on affines (Cohomology of Schemes, Lemma coherent-lemma-quasi-coherent-affine-cohomology-zero) via the higher direct image discussion of Cohomology of Spaces
\item surjectivity of H\textasciicircum{}0 on affines for surjective maps of quasi-coherent sheaves
\item the derived limit criterion for systems of sheaves with surjective transitions and vanishing higher cohomology on a basis (Cohomology on Sites, Lemma sites-cohomology-lemma-derived-limit-suitable-system)
\item the identification of ordinary and derived limits under these conditions (Cohomology on Sites, Lemma sites-cohomology-lemma-inverse-limit-is-derived-limit)
\end{itemize}\medskip
\ifdefined\hypertarget\hypertarget{sourcepointer-1}{}\fi
\noindent\textit{Source reference.} Original lemma 01XB: sources/stacks-chow/tags.txt, source lines 2512--2512. Prerequisite pointer invoked at proof line(s) 608; it names a standard fact of a sibling chapter that the counted argument uses, and every step of the counted proof built on it is retained verbatim in the delivered body.\par\medskip
\ifdefined\hypertarget\hypertarget{sourcepointer-2}{}\fi
\noindent\textit{Source reference.} Original lemma 0A09: sources/stacks-chow/tags.txt, source lines 12627--12627. Prerequisite pointer invoked at proof line(s) 612; it names a standard fact of a sibling chapter that the counted argument uses, and every step of the counted proof built on it is retained verbatim in the delivered body.\par\medskip
\ifdefined\hypertarget\hypertarget{sourcepointer-3}{}\fi
\noindent\textit{Source reference.} Original lemma 0BKY: sources/stacks-chow/tags.txt, source lines 14576--14576. Prerequisite pointer invoked at proof line(s) 616; it names a standard fact of a sibling chapter that the counted argument uses, and every step of the counted proof built on it is retained verbatim in the delivered body.\par\medskip
\ifdefined\hypertarget\hypertarget{sourcepointer-4}{}\fi
\noindent\textit{Source reference.} Original section 071Y: sources/stacks-chow/tags.txt, source lines 9007--9007. Prerequisite pointer invoked at proof line(s) 605; it names a standard fact of a sibling chapter that the counted argument uses, and every step of the counted proof built on it is retained verbatim in the delivered body.\par\medskip
\begin{lemma}
\label{lemma-Rlim-quasi-coherent}
Let $S$ be a scheme. Let $X$ be an algebraic space over $S$.
Let $(K_n)$ be an inverse system of
$D_\QCoh(\mathcal{O}_X)$ with derived limit
$K = R\lim K_n$ in $D(\mathcal{O}_X)$. Assume $H^q(K_{n + 1}) \to H^q(K_n)$
is surjective for all $q \in \mathbf{Z}$ and $n \geq 1$.
Then
\begin{enumerate}
\item $H^q(K) = \lim H^q(K_n)$,
\item $R\lim H^q(K_n) = \lim H^q(K_n)$, and
\item for every affine open $U \subset X$ we have
$H^p(U, \lim H^q(K_n)) = 0$ for $p > 0$.
\end{enumerate}
\end{lemma}

\begin{proof}
Let $\mathcal{B} \subset \Ob(X_\etale)$ be the set of affine objects.
Since $H^q(K_n)$ is quasi-coherent we have $H^p(U, H^q(K_n)) = 0$
for $U \in \mathcal{B}$ by the discussion in
Cohomology of Spaces, Section
\ref{spaces-cohomology-section-higher-direct-image}
and
Cohomology of Schemes, Lemma
\ref{coherent-lemma-quasi-coherent-affine-cohomology-zero}.
Moreover, the maps $H^0(U, H^q(K_{n + 1})) \to H^0(U, H^q(K_n))$
are surjective for $U \in \mathcal{B}$ by similar reasoning.
Part (1) follows from Cohomology on Sites, Lemma
\ref{sites-cohomology-lemma-derived-limit-suitable-system}
whose conditions we have just verified.
Parts (2) and (3) follow from
Cohomology on Sites, Lemma
\ref{sites-cohomology-lemma-inverse-limit-is-derived-limit}.
\end{proof}

\end{document}
